\caption{Commutator spectral separation and finite-data recovery in controlled non-normal systems.
(a) Exact separation $g_C$ between the extreme and interior eigenvalues of the symmetric commutator $C(F)=FF^{\mathsf T}-F^{\mathsf T}F$ as the competing-block commutator strength $c$ approaches that of the planted block. The four plotted strengths are discrete design settings (connected only as a visual guide).
(b) Median principal-angle error of M3 relative to the independently planted plane, $\theta_{\mathrm{M3}}=\theta_{\max}(Q_{\mathrm{M3}},Q_\star)$, for a well-separated ($c=0$) and nearly tied ($c=0.98$) commutator pair across six training regimes. Symbols are medians and whiskers are the archived 95\% independent-base-subject bootstrap intervals. The abscissa is the number of training transition pairs $n=M(T-1)$; the regimes also differ in the number and duration of trajectories, so $n$ alone does not identify the estimation protocol.
(c) Median planted-plane angle for M2 (orange) and M3 (green) in the high-information regime $M=16$, $T=240$. Whiskers are exploratory 95\% percentile intervals, recomputed from frozen rows using 5,000 resamples stratified across $N=16,32,64$ (six independently seeded systems in each dimension; 18 systems per strength).
(d) Number of systems out of 18 for which the conservative commutator perturbation bound on $\sin\theta_{\max}$ is informative (strictly less than one) at two sampling regimes. The bound holds on all 432 fits but is informative on only 36; its vacuity is not a failure of the inequality or proof of poor empirical recovery.
There are 18 independent base systems tested at four paired competition levels and six data regimes ($432$ fitted cases). Plane angles include cases with inadmissible $R_A$; they are not replaced by zero or excluded. The perturbation theorem applies to M3 through the symmetric commutator and is not a general theorem for the nonconvex M2 extractor.}
